	\section{Introducción}
	La inclusión de personas sordas en entornos sociales, educativos y laborales continúa siendo un
	desafío importante en México. La Lengua de Señas Mexicana (LSM), reconocida oficialmente
	como lengua nacional, representa el principal medio de comunicación para esta comunidad. Sin
	embargo, su uso aún no está generalizado entre personas oyentes, lo cual genera barreras
	significativas para la interacción cotidiana. La falta de intérpretes certificados y de herramientas
	tecnológicas accesibles ha motivado la búsqueda de soluciones automatizadas que permitan
	interpretar la LSM de manera eficiente.
	
	De acuerdo con el Instituto Nacional de Estadística y Geografía (INEGI) en 2024,
	aproximadamente el 7.2\% de la población mexicana de 5 años y más vive con alguna
	discapacidad, lo que representa más de 8.8 millones de personas. Entre ellas, el 19.5\% tiene dificultades auditivas y el 9.2\% presenta dificultades para hablar o comunicarse verbalmente, lo
	cual representa una barrera significativa para participar plenamente en la vida cotidiana. La falta
	de intérpretes certificados en Lengua de Señas Mexicana (LSM) y de herramientas accesibles
	para la comunicación agrava esta situación. Esta brecha de accesibilidad evidencia la urgencia de
	desarrollar soluciones tecnológicas que promuevan la comunicación inclusiva sin depender
	exclusivamente de intermediarios humanos.
	
	La Lengua de Señas Mexicana (LSM) es un idioma visual completo, con una estructura
	lingüística, gramatical y cultural propia, desarrollada de forma independiente dentro del contexto
	mexicano. A diferencia del español hablado o escrito, la LSM emplea expresión espacial, gestos
	faciales y movimiento corporal para transmitir significado, sin seguir una traducción literal
	palabra por palabra. Aunque existe el alfabeto dactilológico —útil para deletrear palabras como
	m-a-m-a—, también hay señas específicas para conceptos completos como “mamá” o “comer”.
	Esta lengua no es universal ni mutuamente comprensible con otras como la ASL
	(estadounidense), ya que las señas, expresiones y reglas sintácticas cambian significativamente.
	Incluso en México, puede presentar variaciones regionales, lo que refuerza la necesidad de
	desarrollar soluciones tecnológicas entrenadas específicamente para la LSM, y no adaptar
	sistemas basados en otras lenguas de señas.
	
	La LSM está conformada por diversos elementos lingüísticos, entre los que se incluyen el
	alfabeto dactilológico, las señas léxicas (que representan conceptos completos), las estructuras
	gramaticales propias y los elementos no manuales, como expresiones faciales, mirada y postura
	corporal. Estas características permiten expresar ideas complejas de forma visual y estructurada.
	Además, como cualquier lengua natural, la LSM posee sus propios mecanismos de tiempo verbal,
	negación, preguntas y adjetivación, organizados espacialmente en lugar de linealmente como en
	el español. Comprender esta riqueza estructural es esencial para el desarrollo de sistemas que
	aspiren a interpretarla correctamente.
	
	En los últimos años, diversos enfoques han sido aplicados al reconocimiento del lenguaje de
	señas, desde sensores físicos como guantes con acelerómetros y dispositivos Kinect, hasta
	sistemas basados en video y procesamiento de secuencias gestuales. Si bien estos métodos han
	demostrado efectividad, presentan importantes desventajas: son costosos, difíciles de transportar,
	requieren condiciones de iluminación o espacio controladas, y en muchos casos, demandan
	hardware adicional que limita su accesibilidad. Además, una gran parte de los desarrollos
	existentes se enfocan en la American Sign Language (ASL), dejando de lado la LSM y sus
	particularidades lingüísticas y culturales.
	
	Gracias al avance en el campo de la visión por computadora y el aprendizaje profundo (deep
	learning, DL), se ha abierto una nueva posibilidad para abordar este problema mediante el uso de
	modelos entrenados para procesar imágenes. En particular, las redes neuronales convolucionales
	(CNN) se han consolidado como una de las herramientas más eficaces para la clasificación de
	imágenes, ya que pueden detectar patrones visuales complejos como bordes, formas, texturas y
	posiciones relativas entre elementos.
	

